\documentclass[11pt]{article}
% _preamble.tex --- reglages communs IEEE pour tous les TP du parcours "Improve"
% A inclure via % _preamble.tex --- reglages communs IEEE pour tous les TP du parcours "Improve"
% A inclure via % _preamble.tex --- reglages communs IEEE pour tous les TP du parcours "Improve"
% A inclure via \input{_preamble} JUSTE APRES \documentclass[10pt,journal]{IEEEtran}
% Compilation : tectonic (telecharge les packages automatiquement).
\usepackage[T1]{fontenc}
\usepackage[utf8]{inputenc}
\usepackage[french]{babel}
\usepackage{amsmath,amssymb}
\usepackage{newtxtext,newtxmath}     % rendu Times (proche de Tinos) ; APRES amssymb
\usepackage{graphicx}
\usepackage{booktabs}
\usepackage{array}
\usepackage{xcolor}
\usepackage{listings}
\usepackage{enumitem}
\usepackage{tikz}
\usepackage{pgfplots}
\pgfplotsset{compat=1.18}
\usetikzlibrary{arrows.meta,patterns,calc,decorations.pathmorphing}
\usepackage{cite}
\usepackage[hidelinks]{hyperref}

% --- Noms francais des flottants ---
\renewcommand{\tablename}{Tableau}
\renewcommand{\figurename}{Fig.}
\renewcommand{\lstlistingname}{Listing}
\addto\captionsfrench{\renewcommand{\refname}{Références}}

% --- En-tete de revue commun a tous les TP ---
\newcommand{\improvejournal}{IEEE Transactions on Computational Mechanics and Machine Learning, Vol.~1, No.~1, Juin~2026}

% --- Style code Python (mono, sobre facon IEEE) ---
\definecolor{kw}{HTML}{0033B3}
\definecolor{cm}{HTML}{7A7A7A}
\definecolor{st}{HTML}{067D17}
\lstdefinestyle{py}{%
  language=Python, basicstyle=\ttfamily\footnotesize,
  keywordstyle=\color{kw}, commentstyle=\color{cm}\itshape, stringstyle=\color{st},
  numbers=none, showstringspaces=false, breaklines=true, columns=fullflexible,
  aboveskip=3pt, belowskip=1pt, xleftmargin=2pt,
  literate={é}{{\'e}}1 {è}{{\`e}}1 {à}{{\`a}}1 {ê}{{\^e}}1 {ô}{{\^o}}1 {î}{{\^i}}1 {ç}{{\c c}}1 {É}{{\'E}}1 {_}{{\_}}1}
\lstset{style=py}
 JUSTE APRES \documentclass[10pt,journal]{IEEEtran}
% Compilation : tectonic (telecharge les packages automatiquement).
\usepackage[T1]{fontenc}
\usepackage[utf8]{inputenc}
\usepackage[french]{babel}
\usepackage{amsmath,amssymb}
\usepackage{newtxtext,newtxmath}     % rendu Times (proche de Tinos) ; APRES amssymb
\usepackage{graphicx}
\usepackage{booktabs}
\usepackage{array}
\usepackage{xcolor}
\usepackage{listings}
\usepackage{enumitem}
\usepackage{tikz}
\usepackage{pgfplots}
\pgfplotsset{compat=1.18}
\usetikzlibrary{arrows.meta,patterns,calc,decorations.pathmorphing}
\usepackage{cite}
\usepackage[hidelinks]{hyperref}

% --- Noms francais des flottants ---
\renewcommand{\tablename}{Tableau}
\renewcommand{\figurename}{Fig.}
\renewcommand{\lstlistingname}{Listing}
\addto\captionsfrench{\renewcommand{\refname}{Références}}

% --- En-tete de revue commun a tous les TP ---
\newcommand{\improvejournal}{IEEE Transactions on Computational Mechanics and Machine Learning, Vol.~1, No.~1, Juin~2026}

% --- Style code Python (mono, sobre facon IEEE) ---
\definecolor{kw}{HTML}{0033B3}
\definecolor{cm}{HTML}{7A7A7A}
\definecolor{st}{HTML}{067D17}
\lstdefinestyle{py}{%
  language=Python, basicstyle=\ttfamily\footnotesize,
  keywordstyle=\color{kw}, commentstyle=\color{cm}\itshape, stringstyle=\color{st},
  numbers=none, showstringspaces=false, breaklines=true, columns=fullflexible,
  aboveskip=3pt, belowskip=1pt, xleftmargin=2pt,
  literate={é}{{\'e}}1 {è}{{\`e}}1 {à}{{\`a}}1 {ê}{{\^e}}1 {ô}{{\^o}}1 {î}{{\^i}}1 {ç}{{\c c}}1 {É}{{\'E}}1 {_}{{\_}}1}
\lstset{style=py}
 JUSTE APRES \documentclass[10pt,journal]{IEEEtran}
% Compilation : tectonic (telecharge les packages automatiquement).
\usepackage[T1]{fontenc}
\usepackage[utf8]{inputenc}
\usepackage[french]{babel}
\usepackage{amsmath,amssymb}
\usepackage{newtxtext,newtxmath}     % rendu Times (proche de Tinos) ; APRES amssymb
\usepackage{graphicx}
\usepackage{booktabs}
\usepackage{array}
\usepackage{xcolor}
\usepackage{listings}
\usepackage{enumitem}
\usepackage{tikz}
\usepackage{pgfplots}
\pgfplotsset{compat=1.18}
\usetikzlibrary{arrows.meta,patterns,calc,decorations.pathmorphing}
\usepackage{cite}
\usepackage[hidelinks]{hyperref}

% --- Noms francais des flottants ---
\renewcommand{\tablename}{Tableau}
\renewcommand{\figurename}{Fig.}
\renewcommand{\lstlistingname}{Listing}
\addto\captionsfrench{\renewcommand{\refname}{Références}}

% --- En-tete de revue commun a tous les TP ---
\newcommand{\improvejournal}{IEEE Transactions on Computational Mechanics and Machine Learning, Vol.~1, No.~1, Juin~2026}

% --- Style code Python (mono, sobre facon IEEE) ---
\definecolor{kw}{HTML}{0033B3}
\definecolor{cm}{HTML}{7A7A7A}
\definecolor{st}{HTML}{067D17}
\lstdefinestyle{py}{%
  language=Python, basicstyle=\ttfamily\footnotesize,
  keywordstyle=\color{kw}, commentstyle=\color{cm}\itshape, stringstyle=\color{st},
  numbers=none, showstringspaces=false, breaklines=true, columns=fullflexible,
  aboveskip=3pt, belowskip=1pt, xleftmargin=2pt,
  literate={é}{{\'e}}1 {è}{{\`e}}1 {à}{{\`a}}1 {ê}{{\^e}}1 {ô}{{\^o}}1 {î}{{\^i}}1 {ç}{{\c c}}1 {É}{{\'E}}1 {_}{{\_}}1}
\lstset{style=py}

\begin{document}

\tptitre{1}{1.1}{Flexion d'une poutre : la vérité-terrain analytique}{2 h 30}

\noindent Avant de simuler quoi que ce soit --- et \emph{a fortiori} avant de laisser une IA prédire
à ta place --- il faut une référence dont tu es certain. Ce TP construit cette référence : la
solution analytique d'une poutre en flexion. C'est ta \textbf{vérité-terrain} pour tout le reste
du parcours.

\begin{objectifs}
\begin{itemize}[leftmargin=1.4em,nosep]
  \item Maîtriser la théorie des poutres d'\textbf{Euler-Bernoulli} : flèche, moment, contrainte.
  \item Énoncer ses \textbf{hypothèses} et donc son \textbf{domaine de validité}.
  \item Retrouver à la main la flèche en bout et la contrainte maximale d'une poutre encastrée.
  \item Vérifier ton calcul par un petit script Python.
  \item Découvrir, dès maintenant, l'idée centrale du parcours : \emph{un modèle a un domaine
        au-delà duquel il ment} (ici, Euler-Bernoulli vs poutres courtes).
\end{itemize}
\end{objectifs}

\begin{prereq}
\textbf{Connaissances} : aire, moment quadratique, contrainte normale ($\sigma=F/A$).\\
\textbf{Logiciel} : Python 3 avec \texttt{numpy} (déjà installé). Aucun solveur EF ici --- que de
l'analytique.\\
\textbf{Temps} : $\approx$ 2 h 30. Fais les parties dans l'ordre.
\end{prereq}

\section{Contexte : pourquoi cette poutre est stratégique}
Chez CAE comme dans tout calcul de structure, la question n'est jamais \og{}quel est le résultat ?\fg{}
mais \og{}\emph{ai-je le droit de croire ce résultat ?}\fg{}. La poutre encastrée chargée en bout
(\emph{cantilever}) est le cas-école parfait : sa solution exacte est connue, donc elle te servira
de juge pour tes futurs solveurs EF et tes futurs modèles ML. Si ton code --- ou une IA --- ne
reproduit pas ce cas, tu sais qu'il a tort.

\section{Rappel théorique}

\begin{rappel}[: théorie d'Euler-Bernoulli]
Pour une poutre élancée en flexion, on relie la \textbf{flèche} $w(x)$ au \textbf{moment
fléchissant} $M(x)$ par :
\[
EI\,\frac{d^2 w}{dx^2} = M(x),
\]
où $E$ est le module d'Young (Pa) et $I$ le moment quadratique de la section (m\textsuperscript{4}).
La \textbf{contrainte normale} de flexion à une fibre située à la distance $y$ de l'axe neutre :
\[
\boxed{\;\sigma(x,y) = \frac{M(x)\,y}{I}\;}
\]
La contrainte est maximale là où $|M|$ est maximal et à la fibre la plus éloignée $y=c$ (peau).
\end{rappel}

\subsection{Hypothèses (= le domaine de validité)}
Euler-Bernoulli suppose : (i) petites déformations, (ii) matériau élastique linéaire,
(iii) sections planes restant planes et \emph{perpendiculaires} à la fibre neutre --- donc
\textbf{cisaillement transverse négligé}. Cette dernière hypothèse n'est valable que pour une
poutre \textbf{élancée} : élancement $L/h \gtrsim 10$. En dessous, il faut le modèle de
\textbf{Timoshenko} (qui ajoute le cisaillement) --- sinon Euler-Bernoulli \emph{sous-estime} la
flèche.

\begin{pont}
Retiens cette phrase, c'est tout ton futur métier : \textbf{\og{}Euler-Bernoulli devient faux quand
la poutre est trop courte.\fg{}} C'est exactement la même logique qu'un \emph{surrogate model} qui
ment hors de son domaine d'entraînement (TP 1.5). Un modèle analytique a un domaine ; un modèle ML
aussi. Savoir le tracer, c'est ta valeur ajoutée face à l'IA.
\end{pont}

\subsection{Section rectangulaire}
Pour une section rectangulaire de largeur $b$ et hauteur $h$ :
\[
I = \frac{b\,h^3}{12}, \qquad c = \frac{h}{2}.
\]

\section{Cas d'étude}
Poutre encastrée à gauche (\emph{cantilever}), longueur $L$, chargée par une force ponctuelle $P$
vers le bas à l'extrémité libre.

\begin{center}
\renewcommand{\arraystretch}{1.2}
\begin{tabular}{l l l}
\toprule
Grandeur & Symbole & Valeur \\
\midrule
Longueur & $L$ & $1{,}20$ m \\
Largeur de section & $b$ & $30$ mm \\
Hauteur de section & $h$ & $60$ mm \\
Module d'Young (alu 6061) & $E$ & $69$ GPa \\
Force en bout & $P$ & $800$ N \\
\bottomrule
\end{tabular}
\end{center}

Pour ce chargement, les résultats classiques (à \textbf{démontrer puis appliquer}) sont :
\[
M_{\max} = P L \ \text{(à l'encastrement)}, \qquad
w_{\text{bout}} = \frac{P L^3}{3 E I}, \qquad
\sigma_{\max} = \frac{P L \, (h/2)}{I}.
\]

\section{Partie A --- Calcul à la main}

\begin{exercice}[moment quadratique]
Calcule $I = \dfrac{b h^3}{12}$ pour la section donnée. \textbf{Attention aux unités} : convertis
$b$ et $h$ en mètres avant de calculer. Donne $I$ en m\textsuperscript{4}.
\end{exercice}

\begin{exercice}[moment et contrainte maximale]
\begin{enumerate}[leftmargin=1.4em,nosep]
  \item Donne $M_{\max}=PL$ en N$\cdot$m, et explique \emph{où} il agit sur la poutre.
  \item Calcule $\sigma_{\max}=\dfrac{M_{\max}\,(h/2)}{I}$ en MPa.
  \item La limite d'élasticité de l'alu 6061-T6 est $\approx 240$ MPa. La poutre reste-t-elle
        dans le domaine élastique (hypothèse (ii)) ? Conclus.
\end{enumerate}
\end{exercice}

\begin{exercice}[flèche en bout]
Calcule $w_{\text{bout}} = \dfrac{P L^3}{3 E I}$ en mm. Cette flèche te paraît-elle \og{}petite\fg{}
devant $L$ ? (Vérifie l'hypothèse (i) des petites déformations : $w_{\text{bout}} \ll L$.)
\end{exercice}

\begin{exercice}[le domaine de validité, à la main]
Calcule l'élancement $L/h$ du cas d'étude. Es-tu dans le domaine d'Euler-Bernoulli
($L/h \gtrsim 10$) ? \\
Maintenant imagine une poutre \emph{trapue} : même section mais $L = 0{,}25$ m. Recalcule $L/h$.
Que devrais-tu utiliser à la place d'Euler-Bernoulli, et la flèche réelle serait-elle
\emph{plus grande} ou \emph{plus petite} que celle prédite par Euler-Bernoulli ?
\end{exercice}

\section{Partie B --- Vérification numérique en Python}
Complète le squelette suivant (remplace les \texttt{\#\,TODO}). Le but n'est pas de \og{}coder un
solveur\fg{} mais de \emph{vérifier} tes calculs de la Partie A et d'automatiser l'analyse d'unités.

\begin{lstlisting}
import numpy as np

# --- Donnees du cas d'etude (unites SI) ---
L = 1.20            # m
b = 30e-3           # m
h = 60e-3           # m
E = 69e9            # Pa
P = 800.0           # N
sigma_limite = 240e6  # Pa (limite elastique alu 6061-T6)

# --- TODO 1 : moment quadratique I = b*h^3/12 ---
I = ...

# --- TODO 2 : moment max, contrainte max, fleche en bout ---
M_max = ...
sigma_max = ...
w_bout = ...

# --- Affichage ---
print(f"I        = {I:.3e} m^4")
print(f"M_max    = {M_max:.1f} N.m")
print(f"sigma_max= {sigma_max/1e6:.1f} MPa")
print(f"w_bout   = {w_bout*1000:.2f} mm")

# --- TODO 3 : marge de securite et verification des hypotheses ---
marge = sigma_limite / sigma_max          # > 1 => OK elastique
elancement = L / h
print(f"marge securite = {marge:.2f}")
print(f"elancement L/h = {elancement:.1f}  -> Euler-Bernoulli valide ? {elancement >= 10}")
\end{lstlisting}

\begin{exercice}[confrontation]
Lance le script. Les valeurs Python collent-elles à ton calcul à la main (Partie A) à mieux que
1~\% ? Si non, où est l'erreur --- dans le calcul main ou dans une conversion d'unités ?
\end{exercice}

\begin{exercice}[balayage du domaine]
Modifie le script pour calculer $w_{\text{bout}}$ et $L/h$ pour $L \in \{0{,}2; 0{,}4; 0{,}8;
1{,}2; 1{,}6\}$ m (boucle \texttt{for}). Trace mentalement (ou avec \texttt{matplotlib}) la
frontière $L/h = 10$ : pour quelles longueurs Euler-Bernoulli devient-il douteux ? Tu viens de
tracer ton premier \textbf{domaine de validité}.
\end{exercice}

\begin{piege}
L'erreur n°1 des débutants : mélanger mm et m. $I$ varie comme $h^3$ : un facteur 1000 sur $h$
devient $10^9$ sur $I$. \textbf{Convertis tout en SI dès la première ligne}, jamais à la fin.
\end{piege}

\section{Partie C --- Synthèse}
\begin{exercice}[note d'ingénieur]
Rédige 5 lignes maximum répondant à : \og{}La poutre tient-elle ? À quelle condition ai-je le droit
de croire ce calcul ?\fg{} Mentionne explicitement la marge de sécurité \emph{et} le domaine de
validité (élancement, élasticité, petites déformations).
\end{exercice}

\begin{livrable}
\begin{enumerate}[leftmargin=1.4em,nosep]
  \item Tes calculs à la main des exercices 1 à 4 (photo/scan ou recopie propre).
  \item Le script Python complété et son affichage.
  \item La \og{}note d'ingénieur\fg{} de l'exercice 7 (5 lignes).
  \item Une phrase : \emph{en quoi cette poutre va-t-elle servir de juge à mon futur solveur EF ?}
\end{enumerate}
\end{livrable}

\section{Auto-évaluation}
\begin{center}
\renewcommand{\arraystretch}{1.25}
\begin{tabular}{>{\raggedright\arraybackslash}p{10.5cm} c}
\toprule
Critère & Acquis ? \\
\midrule
J'énonce les 3 hypothèses d'Euler-Bernoulli & $\square$ \\
Je sais dire quand le modèle devient faux (poutre courte $\to$ Timoshenko) & $\square$ \\
J'ai retrouvé $I$, $\sigma_{\max}$ et $w_{\text{bout}}$ à la main & $\square$ \\
Mon script reproduit ces valeurs à $<1\%$ & $\square$ \\
Je comprends le parallèle \og{}domaine d'un modèle analytique = domaine d'un surrogate\fg{} & $\square$ \\
\bottomrule
\end{tabular}
\end{center}

\vfill
\begin{center}\small\itshape
Prochain TP : \textbf{1.2 --- Premier solveur EF 1D (barre en traction)} : on passe de la formule
exacte à la \emph{méthode} qui marche sur n'importe quelle géométrie.
\end{center}

\end{document}
